\documentclass[a5paper,8pt,twoside]{book}
\usepackage{geometry}
\usepackage{csvsimple}
\usepackage{graphicx}
\usepackage[latin9]{inputenc}
\geometry{a5paper, inner=2cm, outer=1.5cm, top=1cm, bottom=1cm}
\usepackage{xstring}
\usepackage[T1]{fontenc}
\usepackage{charter}
\usepackage{longtable}
\usepackage{adjustbox}
\usepackage{ifthen}
\usepackage{titleps}
\usepackage{marginnote}
\usepackage{xcolor}
\setlength{\parindent}{0pt}

%page numbering 
\usepackage{fancyhdr} %page numbering 
\pagestyle{fancy}
\fancyhf{}
\fancyfoot[LE,RO]{\vspace{-0.5cm}\thepage}
\renewcommand{\headrulewidth}{0pt}

%\fancyhead[LE,RO]{\Location} %Add location 

% Create new commands for table of contents
\newcommand*{\nsection}[2]{
	\section*{#1}
	\StrSubstitute{#2}{ }{_}[\Temp]
	\label{\Temp}
	\addcontentsline{toc}{section}{#1}
}

% Create new command to get page reference for each species for TOC
\newcommand*{\pagerefname}[1]{
	 \StrSubstitute{#1}{ }{_}[\TempName]
	 \pageref{\TempName} 
 }

% Create new command to handle empty cells
\newcommand{\ifnotempty}[3]{%
    \IfSubStr{#1}{#2}{#3}{}%
}

% Create new default height and width
\newcommand{\defaultheight}{\0.5\textheight}


% Document begins here
\begin{document}
	% titlepages.tex

\begin{titlepage}
	
%	\begin{figure}[b!]
%		\centering
%		\centerline{\includegraphics[width=0.75\linewidth]{WIP}}
%	\end{figure}

	\includegraphics[width=0.25\linewidth]{Logos/dist_eco} \qquad
	\includegraphics[width=0.25\linewidth]{Logos/biogeo} \qquad
	\includegraphics[width=0.25\linewidth]{Logos/Uni_Bayreuth}
	\vspace{1.5cm}
		
	\csvreader[no head, head to column names, /csv/separator=semicolon]{booklet_title.csv}{1=\Property,2=\Value}{
		\centering
		\IfEq{\Property}{Title}{\Huge\bfseries \Value}{}\\
		\vspace{0.5 cm}
		\IfEq{\Property}{Location}{\Huge \emph{\Value}}{}\\
        }		

\end{titlepage}

\begin{titlepage}
		\csvreader[no head, head to column names, /csv/separator=semicolon]{booklet_title.csv}{1=\Property,2=\Value}{
		\centering
		\IfEq{\Property}{Title}{\Large\Value}{}\\
		\vspace{0.5 cm}
		\IfEq{\Property}{Location}{\large \textbf{\Value}}{}\\
		\IfEq{\Property}{Version}{\large Version \Value}{}\\
		\IfEq{\Property}{Author}{Authors: \large \Value}{}\\
		\IfEq{\Property}{Responsible}{Responsible: \large \Value}{}\\
		\IfEq{\Property}{Contact}{Contact: \large \Value}{}\\
		\IfEq{\Property}{Date}{\large \Value\par\vfill}{}
	}
\end{titlepage}

	\newpage
	
	\section*{Species organised by plant family}

	% Table of contents
    \begin{longtable}{| p{.4\textwidth} | p{.18\textwidth} | p{.25\textwidth} | p{.05\textwidth} |} 
		\hline
		\footnotesize \bfseries Species &
        \footnotesize \bfseries Family &
        \footnotesize \bfseries Local name &
        \footnotesize \bfseries Page \\
		\hline
		\csvreader[head to column names, /csv/separator=semicolon]{species_list.csv}{}{
			\footnotesize{\Species} & 
            \footnotesize{\Family} & 
            \footnotesize{\LocalName} & \footnotesize{\pagerefname{\Species}} \\ \hline
		}
	\end{longtable}

	\newpage
	
	% Iterate over species_list.csv 
	\csvreader[head to column names, /csv/separator=semicolon]{species_list.csv}{}{
	
    \begin{adjustbox}{minipage=[t][0.48\textheight][t]{\textwidth}}
    	       
        \nsection{\emph{\Species} (\Family) \footnotesize{\emph{\Author}}}{\Species}
        \raisebox{\dimexpr\topskip-\height}{\marginnote{\textcolor{\LocationColour}{\textbf{\Location}}}}
         
            \vspace{-0.7cm}
            \ifnotempty{\Synonyms}{}{\footnotesize{\emph{\Synonyms}}\\}
            \footnotesize{\Comments}\\       
            
        \vfill
    \begin{adjustbox}{minipage=[t][\dimexpr\textheight-0.48\textheight-em][t]{\textwidth}}
    	\begin{minipage}{0.5\textwidth}
            \includegraphics[height=\dimexpr\textheight-\ht\strutbox\relax, width=0.9\textwidth, keepaspectratio]{Images/\FileNamePhoto}
            %\includegraphics[height=\fig1width\textheight]{Images/\FileNamePhoto}
        \end{minipage}
       	 \hfill
        \begin{minipage}{0.5\textwidth}
            %\Comments\\
            \includegraphics[height=\dimexpr\textheight-\ht\strutbox\relax, width=0.9\textwidth, keepaspectratio]{Images/\FileNameID}
        \end{minipage}
    \end{adjustbox}
   	 \vfill
    \end{adjustbox}
    \vspace{0.1cm}
}

% Iterate over special_pages.csv

\csvreader[head to column names, /csv/separator=semicolon]{special_pages.csv}{}{
    \begin{adjustbox}{minipage=[t][0.48\textheight][t]{\textwidth}}
        \begin{minipage}{\textwidth}
        \centering
            \includegraphics[
                height=\MaxHeight\defaultheight,
                width=\MaxWidth\textwidth,
                keepaspectratio
            ]{Images/\FileName}
        \end{minipage}
         \hfill
    \end{adjustbox}
    \vspace{0.1cm}
}

\end{document}